\section{MNN 実装計画}

最終実装は MNN CPU inference を使う。offline training は PyTorch、TensorFlow、
ONNX などを使ってよいが、論文で実装可能性を主張するには以下が必要である。

\begin{enumerate}
  \item 入力 tensor shape、10 秒統計スキーマ、系列長 $N$ を固定する。
  \item benign training split だけから normalization constant を凍結する。
  \item 候補モデルを ONNX 等へ export し、MNN converter で固定 shape に変換する。
  \item FP32 と、可能なら Int8/FP16 の MNN model を生成する。
  \item 同一入力で offline framework と MNN の anomaly score を比較し、許容誤差を定義する。
  \item model file、weight、tensor、workspace、I/O buffer、score history を含む
        model-owned peak memory を測定する。
  \item peak が 500 KB を超える場合、feature 数、系列長 $N$、hidden size、
        architecture を縮小して再評価する。
\end{enumerate}

実装時の model interface は全候補で共通にする。PyTorch 側の \texttt{forward} は
正規化済み 10 秒統計列
$X \in \mathbb{R}^{B_{\mathrm{batch}} \times N \times D}$ を受け取り、
同じ shape の $\hat{X}$ だけを返す。anomaly score、しきい値判定、
channel-wise error 集計は MNN graph の外で実装する。feature slice は固定順序とし、
count/bytes、read/write ratio、mean LBA、mean length、窓間差分、
optional compression/padding を同じ index に保つ。ONNX/MNN export は
batch size 1、固定 $N$、固定 $D$ とし、動的軸を使わない。metadata-only 実験では
optional slot を zero-fill し、対応する loss weight と score weight を 0 にする。

この手順では、MNN ランタイムのコードサイズや共通 heap は除外する。
一方で、モデルファイルが flash 上に常駐し RAM に展開されない構成か、
ロード時に RAM コピーされる構成かは target device に依存するため、
実装論文では両者を明記する。
